\subsection{Microcontroller Selection}
Four properties of the ATmega32 map onto the requirements of Section~\ref{sec:intro}. Its 10-bit ADC
is multiplexed across eight channels, removing the need for an external converter. Its 32
general-purpose I/O lines are enough, but only just, to host a 32-key matrix alongside two displays
and a control button. Its hardware TWI peripheral is what lets the doctor-side display leave the
parallel bus. And it runs from a factory-calibrated internal RC oscillator, so no crystal is
required~\cite{microchip2020}.

\subsection{I/O and Resource Allocation}
Table~\ref{tab:pinout} gives the complete allocation. Twenty-three of the 32 available lines are in
active use and two more are reserved for the UART, leaving seven free.

\begin{table}[t]
  \caption{ATmega32 I/O allocation: 23 lines active, two reserved, seven free.}
  \label{tab:pinout}
  \small
  \begin{tabular}{@{}llcl@{}}
    \toprule
    Port lines & Function & Lines & Direction \\
    \midrule
    PB0--PB7      & Keypad column drive        & 8 & Output \\
    PA4--PA7      & Keypad row sense           & 4 & Input (pull-down) \\
    PA0 (ADC0)    & LM35 temperature           & 1 & Analog in \\
    PA1 (ADC1)    & XD-58C pulse sensor        & 1 & Analog in \\
    PA2           & Doctor's call button       & 1 & Input (active low) \\
    PC0--PC1      & I\textsuperscript{2}C, doctor LCD & 2 & Bidirectional \\
    PD2--PD7      & Patient LCD, 4-bit parallel & 6 & Output \\
    \midrule
    PD0--PD1      & UART (RXD/TXD), reserved   & 2 & Reserved \\
    PA3, PC2--PC7 & Unallocated                & 7 & n/a \\
    \bottomrule
  \end{tabular}
\end{table}


The keypad dominates the budget, taking twelve lines, more than a third of the total, to drive eight
column outputs and sense four row inputs. That is the price of direct alphanumeric entry, and
Section~\ref{sec:firmware} explains what it buys. Port A is heavily overloaded, carrying both analog
inputs, the four row sense lines and the doctor's button. This is why the keypad driver masks its
direction register rather than writing it wholesale, since an unguarded write to \texttt{DDRA} would
reconfigure the ADC pins mid-acquisition. The display interface asymmetry also pays here. Had the
doctor-side module used the parallel interface rather than an I\textsuperscript{2}C backpack, active
usage would rise from 23 lines to 27, or 29 counting the reserved UART pair, leaving three of 32 free
and no room for an expansion header.

\subsection{Analog Front End}
The LM35 produces 10~mV per degree Celsius with no external circuitry, so its output connects
directly to ADC channel~0~\cite{ti2022}. The XD-58C is an integrated photoplethysmographic module
with on-board amplification and filtering, and its analog output connects to channel~1. The converter
uses the internal 2.56~V bandgap reference, which sets the least significant bit to 2.5~mV and,
against the LM35 transfer function, gives a quantisation step of exactly $0.25^{\circ}$C. AREF is
decoupled to ground through a 1~$\mu$F capacitor as the datasheet specifies~\cite{microchip2020}.

One configuration choice deserves to be flagged rather than defended. The ADC prescaler is left at
its minimum division factor of two, which at 8~MHz yields a 4~MHz conversion clock, whereas the
datasheet specifies 50--200~kHz for full 10-bit accuracy~\cite{microchip2020}. The converter is
therefore clocked well outside the range in which its rated resolution holds. We identify this as a
likely contributor to the temperature spread reported in Section~\ref{sec:results} and as the
cheapest available accuracy improvement, since it costs three bits in one register and no hardware.

\subsection{Clock Configuration and Baud-Rate Error}
\label{sec:calibration}
The internal oscillator ships selected at 1~MHz, and we reconfigured it to 8~MHz through the CKSEL
fuse bits. The reasoning is reproducible from the datasheet, so we set it out here.

\begin{table}[t]
  \caption{UART divisor error at the factory 1~MHz clock versus the fuse-selected
    8~MHz clock. Rows marked $\dagger$ exceed the datasheet's recommended
    $\pm 2.0\%$ receiver error tolerance~\cite{microchip2020}.}
  \label{tab:baud}
  \small
  \begin{tabular}{@{}llcrrr@{}}
    \toprule
    $f_{\mathrm{osc}}$ & Target & U2X & UBRR & Actual & Error \\
    \midrule
    1~MHz & 38400 & 0 & 1  & 31250 & $-18.6\%^{\dagger}$ \\
    1~MHz & 38400 & 1 & 2  & 41667 & $+8.5\%^{\dagger}$ \\
    \midrule
    8~MHz & 38400 & 0 & 12 & 38462 & $+0.16\%$ \\
    8~MHz & 38400 & 1 & 25 & 38462 & $+0.16\%$ \\
    \bottomrule
  \end{tabular}
\end{table}


The ATmega32 derives its UART bit clock by integer division of the system clock. In standard mode the
divisor is $\mathrm{UBRR} = \lfloor f_{\mathrm{osc}}/(16 \times \mathrm{BAUD}) + 0.5 \rfloor - 1$,
with 16 replaced by 8 when the double-speed bit U2X is set~\cite{microchip2020}. Because UBRR is an
integer the achievable rate is quantised, coarsely so when $f_{\mathrm{osc}}$ is low, and the
datasheet's recommended ceiling on receiver error for an 8-bit frame is $\pm 2.0\%$.
Table~\ref{tab:baud} works this through. At 1~MHz, 38400~bps is unreachable, because the nearest
standard-mode divisor yields 31250~bps for an error of $-18.6\%$, roughly nine times the limit, and
double-speed mode improves it only to $+8.5\%$. At 8~MHz the same rate is reachable within
$+0.16\%$ in either mode.

One consequence is worth recording. The internal oscillator is specified to roughly
$\pm 3\%$~\cite{microchip2020}, orders of magnitude looser than the tens of parts per million a quartz
crystal offers. That is adequate for UART framing once the divisor error is this small, but it
propagates into every delay-loop-derived interval in the firmware, including the pulse sampling time
base. The serial path itself is unused by the delivered offline firmware, and establishing that
38400~bps is attainable keeps the optional uplink of Section~\ref{sec:limitations} available without a
hardware revision.
